\chapter{统计描述、过程波动与运营规律}
\label{ch:week05-statistics-dashboard}

描述统计通过分布、位置、离散程度和时间变化刻画经营活动。订单周期的均值与
高分位数分别反映总体时间消耗和长周期订单，产品分组比较区分组内变化与组合
变化，设备利用率和等待时间共同描述运行负荷。

变量分布反映典型水平与离散程度；结构和时间分解区分组合变化、随机波动与过程
信号；存量、流量和等待时间的关系限定运营解释。看板应保留这些区别，并注明
描述性发现及其仍待检验的假设。

\section{经营数据的统计描述}

描述统计以对象含义、计算方法与质量状况明确的观测资料为基础。相关条件见第2至4章。指标、图形、解释与行动的关系为：
\[
\text{可信指标}
\rightarrow
\text{描述统计}
\rightarrow
\text{图表编码}
\rightarrow
\text{经营解释}
\rightarrow
\text{管理行动}.
\]

统计描述概括已观察资料的水平、分布与变化。图形比较可以确定异常集中的位置，并提出后续假设；假设的因果含义由相应的实验或观测性研究检验。

\begin{definition}[描述性管理分析]
描述性管理分析是以可信指标为输入，通过统计摘要、分组比较、时间趋势、异常识别和可视化界面，刻画经营状态并提出后续行动问题的分析活动。
\end{definition}

异常记录和组间差异为进一步研究提供线索。确认干预效果需要第6章和第8章的研究设计，比较具体行动则需要第9章的成本与约束分析。

\section{经验分布与描述统计}
\subsection{经营指标作为观测变量}
\subjectindex{miaoshutongji}{描述统计}

设某指标在 \(n\) 个日期、客户、区域或产品类别上的观测为
\[
x_1,x_2,\ldots,x_n.
\]
这些观测可以来自同一对象的时间序列，也可以来自不同对象在同一时间窗口的横截面。时间序列关注经营节奏和变化方向，横截面关注对象之间的差异。二者不能混用。例如，日收入的波动不能直接解释为区域之间的结构差异；区域收入差异也不能直接说明某区域正在改善或恶化。

\begin{convention}[观测单位]
每个统计摘要都必须声明观测单位。观测单位可以是日期、订单、客户、产品、门店、区域或客户日期组合。若观测单位不清楚，均值、分位数、异常值和趋势都无法被正确解释。
\end{convention}

本章样例数据 \path{data/sample/daily_operations.csv} 的观测单位是一行一个经营日期。它适合讨论日收入、日订单、准时交付率、平均延期、待处理订单和质量问题随时间的变化，但不适合直接回答哪个客户贡献了异常收入。要回答客户问题，需要回到订单事实表和客户维度。

\subsection{经验分布与统计量}
\subjectindex{jingyanfenbu}{经验分布}
对 \(n\) 个观测值 \(x_1,\ldots,x_n\)，经验分布函数定义为
\begin{equation}
F_n(u)=\frac1n\sum_{i=1}^n\mathbf1\{x_i\le u\}.
\label{eq:empirical-cdf}
\end{equation}
它回答“观测对象中有多大比例不超过 \(u\)”的问题。若 \(x_i\) 是订单周期，
\(F_n(5)\) 就是样本中五天以内完成的比例。这个描述本身不要求样本独立同分布；
但把它外推为未来订单的概率分布，需要另行说明样本如何获得、环境是否稳定。

均值、中位数和分位数概括同一分布的不同特征。均值与总成本、总工时相联系，
中位数反映分布中点，高分位数刻画长周期对象。管理问题决定应选用或并列呈现
哪些统计量。

\begin{example}[同一个样本中的不同服务含义]
五张订单的周期为 \(1,1,2,2,14\) 天。平均周期为4天，中位数为2天，五天内完成的
比例为80\%。均值描述总周期20天除以订单数，中位数描述典型位置，比例描述服务
承诺的达成程度。若只展示“典型订单两天完成”，会遗漏第五张订单；若只展示均值，
又可能让人误以为大多数订单接近四天。这些差异应在一张分布图中同时保留。
\end{example}

此外，订单加权和金额加权描述的是不同总体。将大订单赋予更大权重可以刻画资金
暴露，却不能把结果称为“随机抽取一张订单的平均周期”。报告中的每一种权重都应
能回答一个明确问题；无理由地更换权重，与第~\ref{ch:week03-sql-kpi}~章讨论的
更换分母具有相同风险。

\subsection{均值与中位数}
\subjectindex{junzhi}{均值}
\subjectindex{zhongweishu}{中位数}

描述统计首先要回答典型水平。样本均值为
\[
\bar x=\frac{1}{n}\sum_{i=1}^n x_i.
\]
均值适合描述可加指标的平均水平，例如日收入、日订单量或每周质量问题数。它的优点是使用全部观测，并且与总量之间有明确关系：
\[
\sum_{i=1}^n x_i=n\bar x.
\]
均值与总量的对应关系可用于预算、产能和目标差距分析。

但均值对极端值敏感。若某一天出现大客户订单或一次性促销，均值会被明显拉高。
样本中位数 \(Q_2\) 取为排序后的中间值；样本量为偶数时，取中间两个值的平均。
中位数可补充描述典型日或典型客户。若均值明显高于中位数，可以将右偏或少数高值
影响作为待检查的分布特征，但仍应结合实际分布图判断。

\begin{example}[日收入的中心位置]
第5章样例中，总收入为 71350，观测天数为 12，平均日收入为 5945.83，中位数为 5100。均值高于中位数，提示收入分布受到高收入日期影响。这个发现本身不说明经营改善，只说明需要下钻到高收入日期对应的订单、客户或产品。
\end{example}

\subsection{方差、标准差与变异系数}
\subjectindex{fangcha}{方差}
\subjectindex{bianyixishu}{变异系数}

只有中心位置不足以支持管理判断。两个团队的平均收入可能相同，但一个团队稳定，另一个团队高度波动。当 \(n\ge2\) 时，样本方差和标准差定义为
\[
s^2=\frac{1}{n-1}\sum_{i=1}^n(x_i-\bar x)^2,\qquad s=\sqrt{s^2}.
\]
标准差与原指标具有相同单位，用于描述观测值偏离均值的总体幅度。若日收入的标准差为 2987.05，则说明日收入围绕均值有较大波动。

当比较不同量纲或不同水平的指标时，可使用变异系数
\[
CV=\frac{s}{\bar x}.
\]
变异系数描述相对波动，适用于具有真实零点且均值为正的比率尺度变量，例如收入和订单量。
摄氏温度等零点可任意平移的变量不适合直接比较变异系数；当均值接近0或指标可以取负值时，
该比值也会失去稳定解释。

\begin{remark}
标准差刻画统计波动。经营风险还取决于波动发生在什么对象、是否集中在关键客户、是否超出服务承诺、是否触发现金流或库存约束。因此，应结合波动对应的业务对象和约束评价其经营影响。
\end{remark}

\subsection{分位数、IQR 与异常识别}
\subjectindex{fenweishu}{分位数}
\subjectindex{IQR}{IQR（四分位距）}

分位数描述分布结构。设 \(Q_1\)、\(Q_2\)、\(Q_3\) 分别为第一四分位数、中位数和第三四分位数，则四分位距为
\[
IQR=Q_3-Q_1.
\]
IQR 衡量中间 50\% 观测的跨度，对极端值不敏感。常见经验规则将低于
\[
Q_1-1.5IQR
\]
或高于
\[
Q_3+1.5IQR
\]
的观测标记为异常值。

\begin{definition}[描述性异常]
描述性异常是相对于某个经验规则、目标线、历史基线或业务阈值显得不寻常的观测。异常识别确定需要进一步调查的观测，其成因可结合分层数据、业务记录和研究设计分析。
\end{definition}

样本分位数存在不同插值约定。本章配套程序采用上下半样本中位数法：先排序，
偶数样本分成等长两半；奇数样本排除整体中位数后分成两半，再分别取其中位数为
\(Q_1,Q_3\)。不同软件采用的分位数定义可能不同，比较结果时应统一插值约定。
在本章12天样例中，收入的 \(Q_1=4100\)，\(Q_3=6975\)，因此 \(IQR=2875\)。
IQR 规则识别出13500为异常高收入。管理者应进一步查询客户、促销、重复入账和
同期履约情况；异常值标记本身不提供原因。

\section{结构变化、抽样波动与过程信号}
\subsection{比率指标与分母}

看板中常有准时率、毛利率、缺陷率、转化率等比率型指标。比率指标不能简单平均，计算依据见第~\ref{ch:week03-sql-kpi}~章。看板应同时展示比率、分子和分母。若某区域准时率为 \(100\%\)，但只有 1 个交付订单，它的管理含义不同于 300 个订单上的 \(96\%\)。

设分组 \(g\) 的分子为 \(N_g\)，分母为 \(D_g\)。总体比率应为
\[
R=\frac{\sum_g N_g}{\sum_g D_g}.
\]
若直接平均各组比率
\[
\frac{1}{G}\sum_{g=1}^G \frac{N_g}{D_g},
\]
就会把不同分母的组赋予相同权重，从而造成误读。

\begin{table}[htbp]
\centering
\caption{比率指标的看板展示要求}
\begin{tabular}{P{0.24\linewidth}P{0.32\linewidth}P{0.32\linewidth}}
\toprule
\rowcolor{BDMTableHead}
展示元素 & 作用 & 管理风险\\
\midrule

比率值 & 显示当前比例水平。 & 单独展示可能造成小样本误读。\\
分子 & 说明事件发生次数。 & 隐藏分子会掩盖问题规模。\\
分母 & 说明比较基础。 & 分母太小时比较不稳定。\\
目标线 & 提供评价基准。 & 缺少目标线时难以判断达成程度。\\
质量状态 & 说明是否可发布。 & 质量检查未通过时，图表可能传播错误信息。\\
\bottomrule
\end{tabular}
\end{table}

经营报告应说明指标的可加性、分子与分母、样本量是否足以支持比较，以及是否存在未通过质量检查的记录。

\subsection{产品组合与总体指标}
\subjectindex{xinpusen}{辛普森反转}
\subjectindex{chanpinzuhe}{产品组合}
令产品族 \(g\) 在时期 \(t\) 的占比为 \(w_{gt}\)，良率为 \(p_{gt}\)，则总体良率
\[
p_t=\sum_gw_{gt}p_{gt},\qquad \sum_gw_{gt}=1.
\]
在两期相减后加减 \(\sum_gw_{g0}p_{g1}\)，得到精确分解
\[
p_1-p_0
=\underbrace{\sum_gw_{g0}(p_{g1}-p_{g0})}_{\text{固定初期组合的组内变化}}
+\underbrace{\sum_gp_{g1}(w_{g1}-w_{g0})}_{\text{按末期水平计的组合变化}}.
\]
这是一个选定基准的分解，换用另一基准会重新分配交互项；总体差异不变。

\begin{example}[两类产品都改善，总体却下降]
初期简单产品良品为 \(90/100\)，复杂产品为 \(1/10\)，总体为 \(91/110\approx82.73\%\)。末期二者变成 \(19/20\) 与 \(40/200\)，组内良率分别从 \(90\%\) 升至 \(95\%\)、从 \(10\%\) 升至 \(20\%\)，总体却降至 \(59/220\approx26.82\%\)。原因是复杂产品占比大幅增加。该算例中，两类产品的占比变化决定了总体变化的方向。
\end{example}

经营看板应同时显示实际总体指标、固定组合的标准化指标和产品占比。这三项分别描述总体表现、相同组合下的组内变化以及构成变化；干预效果的识别需进一步研究处理分配机制。

\subsection{经营波动与移动平均}
\subjectindex{yidongpingjun}{移动平均}

管理者通常既关心当前水平，也关心趋势。对时间序列 \(x_t\)，简单移动平均可写为
\[
MA_t^{(k)}=\frac{1}{k}\sum_{j=0}^{k-1}x_{t-j}.
\]
移动平均可以平滑短期波动，使趋势较易辨认，但也可能掩盖突发异常。因此，看板应同时展示移动平均和原始观测。

移动平均还有边界与滞后问题。按上式定义的后向窗口在最初 \(k-1\) 个时点无法完整计算；可以保留缺失值，也可以另行约定使用较短窗口并改变分母，二者应在图注中说明。若某天收入突然升高，该高值会进入当天及后续 \(k-1\) 天的移动平均，而不会改变过去已经发布的后向均值。阅读时应同时观察原始观测和移动平均线，避免把平滑曲线误解为原始经营过程。

\begin{convention}[平滑与异常并存]
用于经营监控的时间序列视图应同时保留原始观测、平滑趋势和异常标记。原始观测用于追踪事实，平滑趋势用于识别节奏，异常标记用于触发下钻。
\end{convention}

\subsection{统计过程控制图}
\subjectindex{kongzhitu}{控制图}
经营目标是希望达到的水平，控制界限是一个过程模型下预期波动的范围。二者不能互换。考虑第 \(t\) 天检验 \(n_t\) 件产品，不良件数为 \(D_t\)。在基准过程稳定、各件独立且有共同不良概率 \(p_0\) 的二项模型下，
\[
\hat p_t=D_t/n_t,\qquad
\operatorname{Var}(\hat p_t)=p_0(1-p_0)/n_t.
\]
由此得到常见的近似三标准差控制界限
\[
L_t=\max\{0,p_0-3\sqrt{p_0(1-p_0)/n_t}\},\qquad
U_t=\min\{1,p_0+3\sqrt{p_0(1-p_0)/n_t}\}.
\]
实践中 \(p_0\) 常由一段经过审核的稳定历史估计。NIST/SEMATECH 手册系统讨论了过程监控的背景与条件\cite{nist05}。

\begin{figure}[htbp]
\centering
\begin{tikzpicture}
\begin{axis}[width=0.88\linewidth,height=6cm,xlabel={日序号},
ylabel={不良比例},xmin=0.8,xmax=8.2,ymin=0,ymax=0.12,
xtick={1,...,8},ytick={0,.02,.04,.06,.08,.10,.12},
yticklabels={0,2\%,4\%,6\%,8\%,10\%,12\%},
legend style={at={(0.5,1.03)},anchor=south,legend columns=3}]
\addplot[BDMTeal,mark=*] coordinates {(1,.04)(2,.05)(3,.06)(4,.044)(5,.056)(6,.10)(7,.052)(8,.048)};
\addlegendentry{观测}
\addplot[BDMGray,dashed,domain=1:8] {.05};
\addlegendentry{基准}
\addplot[BDMAmber,dashdotted,domain=1:8] {.05+3*sqrt(.05*.95/500)};
\addlegendentry{控制界限}
\addplot[BDMAmber,dashdotted,domain=1:8,forget plot] {.05-3*sqrt(.05*.95/500)};
\addplot[only marks,mark=*,BDMAmber,mark size=2.5pt,forget plot] coordinates {(6,.10)};
\node[font=\footnotesize,text=BDMAmber,anchor=south east] at (axis cs:5.9,.104) {第6天：50/500};
\end{axis}
\end{tikzpicture}
\caption{教学控制图：每天 500 件、基准不良率 \(5\%\)；第6天值得调查，图本身不识别原因。}
\end{figure}

若每天只抽 20 件，正态近似可能很差；同一批次内缺陷相关会造成过度离散\footnote{过度离散是指数据的实际波动大于所用概率模型规定的方差。例如，同批产品缺陷相关时，不良件数的波动可能超过独立二项模型的预测范围。}；基准本身漂移也会使固定界限失效。此时应考虑精确二项界限、分层或适当的过程模型。反复监控带来累计误报问题，第~\ref{ch:week06-ab-testing}~章将讨论时间一致推断。

\subsection{比较基准、目标线与经营解释}

描述统计需要比较基准。基准限定指标的比较和评价含义。常见基准包括：
\begin{itemize}
  \item 历史基准：与上周、上月、去年同期比较；
  \item 计划基准：与预算、目标、服务水平协议比较；
  \item 同类基准：与区域、门店、产品类别或客户分层比较；
  \item 统计基准：与均值、分位数、控制界限或异常阈值比较。
\end{itemize}

同一个指标在不同基准下可能得到不同解释。收入高于上周并不一定说明达成预算；准时率高于某区域平均也不一定达到客户承诺。看板应明确比较对象，避免不同基准下的判断混用。

\begin{table}[htbp]
\centering
\caption{经营解释中的常见基准}
\begin{tabular}{P{0.22\linewidth}P{0.30\linewidth}P{0.36\linewidth}}
\toprule
\rowcolor{BDMTableHead}
基准 & 适用问题 & 风险\\
\midrule

上期值 & 是否改善或恶化？ & 季节性和工作日结构可能干扰判断。\\
目标线 & 是否达成管理承诺？ & 目标若未更新，可能失去约束力。\\
同类对象 & 哪些区域或产品异常？ & 对象规模、组合和质量状态可能不同。\\
分位数 & 观测是否罕见？ & 罕见不等于错误，也不等于原因。\\
质量检查 & 指标是否可发布？ & 忽略质量状态会造成错误信任。\\
\bottomrule
\end{tabular}
\end{table}

\section{制造系统中的存量、流量与等待}

\subsection{Little 定律的面积证明}
\subjectindex{Little}{Little定律}
把工厂视为一个固定边界内的系统，\(L(t)\) 为时刻 \(t\) 的在制批次数。观察足够长时间 \(T\)，则 \(\int_0^T L(t)\,\mathrm dt\) 是所有批次在窗口内停留时长的总和：同一面积既能沿时间横切，也能按批次竖切。若系统稳定、长期平均存在、边界上未完成批次的贡献除以 \(T\) 趋于零，便得到
\[
\bar L=\lambda\bar W.
\]
其中 \(\lambda\) 是同一边界和单位下的长期有效吞吐率，\(\bar W\) 是平均系统停留时间。这是 Little 定律\cite{little05}，不要求泊松到达或指数服务，但要求口径与长期边界条件一致。

一个稳定工站平均 WIP 为 240 批，每天完成 60 批，则平均周期为 \(240/60=4\) 天。存量、吞吐量与周期必须对应相同工站范围和批次单位，观察期间还应满足稳定平均条件。若处于集中开工或扩产阶段，应先研究在制品随时间的积累过程，再确定适用的统计期间。

\subsection{利用率与等待时间}
\subjectindex{MM1}{M/M/1队列}
\subjectindex{Kingman}{Kingman近似}
\subjectindex{liyonglv}{利用率}
进一步采用 \(M/M/1\) 教学模型：到达为强度 \(\lambda\) 的泊松过程\footnote{齐次泊松过程假定不相交时间区间内的到达次数相互独立，且次数服从均值与区间长度成正比的泊松分布；相邻到达间隔独立并服从同一指数分布。}，单台设备服务时间独立且服从均值 \(1/\mu\) 的指数分布，二者独立，先到先服务，无容量上限，且 \(\rho=\lambda/\mu<1\)。令 \(\pi_n\) 为稳态有 \(n\) 批的概率，相邻状态流量平衡给出 \(\pi_n\lambda=\pi_{n+1}\mu\)，所以 \(\pi_n=\pi_0\rho^n\)。利用几何级数归一化得到 \(\pi_0=1-\rho\)，进而
\[
\bar L=\sum_{n=0}^\infty n(1-\rho)\rho^n
=\frac{\rho}{1-\rho}.
\]
再用 Little 定律计算系统时间，并减去平均服务时间 \(1/\mu\)，分别得到
\[
\bar W=\frac1{\mu-\lambda},\qquad
\bar W_q=\frac{\rho}{\mu-\lambda}.
\]
若 \(\mu=10\) 批/小时，到达率从 8 增加到 9.5，利用率只从 \(80\%\) 增至 \(95\%\)，平均系统时间却从 0.5 小时增至 2 小时。提高局部利用率的收益必须与等待、WIP 和 OTD 的损害共同评估。

\begin{figure}[htbp]
\centering
\begin{tikzpicture}
\begin{axis}[width=0.85\linewidth,height=6cm,xlabel={利用率 \(\rho\)},
ylabel={系统时间 / 平均服务时间},xmin=0,xmax=1,ymin=0,ymax=25,
xtick={0,.2,.4,.6,.8,1},xticklabels={0,20\%,40\%,60\%,80\%,100\%}]
\addplot[BDMTeal,thick,domain=0:.95,samples=100] {1/(1-x)};
\addplot[only marks,BDMAmber] coordinates {(.8,5)(.95,20)};
\draw[BDMRule,dashed] (axis cs:.8,0) -- (axis cs:.8,5) -- (axis cs:0,5);
\draw[BDMRule,dashed] (axis cs:.95,0) -- (axis cs:.95,20) -- (axis cs:0,20);
\node[font=\footnotesize,anchor=south east,fill=white,inner sep=2pt]
  at (axis cs:.78,5.3) {80\%：5倍};
\node[font=\footnotesize,anchor=south east,fill=white,inner sep=2pt]
  at (axis cs:.92,20.3) {95\%：20倍};
\end{axis}
\end{tikzpicture}
\caption{\(M/M/1\) 模型的平均系统时间与利用率。纵轴为等待和服务时间之和相对于平均服务时间的倍数。}
\end{figure}

对到达间隔与服务时间各自独立同分布、两序列相互独立、先到先服务且 \(\rho<1\) 的 \(GI/G/1\) 系统\footnote{排队模型记号中，GI 表示一般分布且相互独立的到达间隔，G 表示一般服务时间分布，1 表示单个服务台。此处其余独立性、服务次序及稳定性条件由正文另外规定。}，在相应有限二阶矩等条件下，Kingman 近似给出平均排队等待时间的估计
\[
\bar W_q\approx
\frac{c_a^2+c_s^2}{2}\frac{\rho}{1-\rho}\E[S],
\]
其中 \(c_a,c_s\) 是到达间隔与服务时间的变异系数。该式是近似而非普遍恒等式，主要用于理解变异性、利用率与服务时间如何共同影响等待。产品混流、故障和换线会使真实产线复杂得多，应通过数据和仿真校准。

\subsection{服务能力的经济分析}
\label{sec:queue-cost}
\subjectindex{paiduimoxing}{排队模型}
在前述稳定 M/M/1 模型中，\(W=1/(\mu-\lambda)\) 是平均系统逗留时间。
保持到达率 \(\lambda\) 不变，其对服务率的导数为
\[
\frac{\mathrm dW}{\mathrm d\mu}=-\frac1{(\mu-\lambda)^2}.
\]
这表明增加相同幅度的能力，在接近拥堵时带来的平均逗留时间改善更大。
同时也说明，不能根据低负荷下的平均周期直接推断高负荷时仍有同样表现。
这里的非线性来自排队累积，不需要假设设备本身变慢。

进一步假定可连续调整服务能力，每单位时间的能力成本为 \(k\mu\)，每件在系统中的
对象每单位时间产生成本 \(h\)。在 \(k,h,\lambda>0\) 且 \(\mu>\lambda\) 时，
利用 \(L=\lambda/(\mu-\lambda)\)，总成本率为
\[
J(\mu)=k\mu+\frac{h\lambda}{\mu-\lambda}.
\]
求导得到 \(J'(\mu)=k-h\lambda/(\mu-\lambda)^2\)，而
\(J''(\mu)=2h\lambda/(\mu-\lambda)^3>0\)，所以唯一内点最优能力为
\begin{equation}
\mu^*=\lambda+\sqrt{\frac{h\lambda}{k}}.
\label{eq:queue-capacity}
\end{equation}

若 \(\lambda=8\)、\(h=1\)、\(k=2\)，则 \(\mu^*=10\)，利用率为80\%，
总成本率为24。把能力压到9，能力成本降低，但总成本率上升到26；把能力提高到12，
总成本率也为26。最优方案保留一定富余能力，使增加能力的边际成本与减少积压的边际收益相等。

式~\eqref{eq:queue-capacity}~适用于连续服务率和线性能力成本。若设备数量只能取
整数，或存在换型、停机与非线性成本，应在相应可行集合上重新比较方案。
第~\ref{ch:week09-optimization}~章将进一步讨论多个资源约束下的选择。

\section{人工智能服务的运营度量}
\label{sec:ai-service-operations}
\subjectindex{TTFT}{TTFT（首输出时间）}
把企业问数系统作为持续运行的服务，可区分请求、完整任务和业务结果三个单位。
一个用户问题可能包含模型生成、数据库查询、文档检索及人工批准等多次调用。
请求到达后到首个输出 token 的时间称为 TTFT（time to first token）\footnote{token 是模型处理文本的基本单位，由所用分词器划分，可能对应一个字、一个词或词的一部分。TTFT 只度量首次输出的等待，不等于完整回答的生成时间。}；从提交到全部回答返回为请求完成时间；
从提出问题到取得通过验收的业务结果为任务完成时间。后两个时间都应记录
失败与重试。李博杰对预填充、逐步生成及任务负载的区分有助于确定测量边界
\cite{aiinfra2026}。

对单次、串行执行的模型调用，设输入 \(p\) 个 token、输出 \(o\) 个 token，
在指定模型和并发条件下观测的预填充与生成速度分别为 \(r_p,r_d\)，可以用
\[
  s\approx t_0+\frac{p}{r_p}+\frac{o}{r_d}
       +t_{\rm query}+t_{\rm tool}
\]
估算无排队时的服务时间。这里 \(t_0\) 是网络与调度的固定耗时；若查询或工具
并行运行，应改用执行依赖图求关键路径。\(r_p\) 与 \(r_d\) 是特定负载下的
观测值，不能把硬件标称算力直接代入。系统还须分别记录首输出、完整任务时间
的分布及第95百分位数；均值相加也不能给出总体的第95百分位数。

设稳定系统每秒收到 \(\lambda\) 个任务，平均完成时间为 \(W\)，则 Little
定律给出平均在途任务数 \(L=\lambda W\)。若仅需估计同时占用服务位置的数量，
可先用无排队服务时间 \(s\) 计算负载 \(\lambda s\)。假设每小时有900个请求、
每次无排队服务时间约6.19秒，则平均同时占用的服务位置约为1.55个。
有四个等效位置时，按总可用服务时间计算的占用比约为0.39。
这个比例不等于任务的第95百分位完成时间；突发到达、服务时间变异、位置共享
及工具等待均可能显著改变尾部延迟。上述占用估计假设查询与工具等待期仍占用
同一服务位置；若会释放，需对模型和工具分别测量占用。

每次请求的费用应按输入、输出、查询与工具调用分别核算。设输入与输出每百万 token 的单价分别为
\(c_p,c_o\)，查询和工具的单位费用为 \(c_q,c_t\)，则
\[
  c_{\rm request}=
    \frac{p c_p+o c_o}{10^6}+c_q+c_t.
\]
若失败任务必须重做，还要把失败尝试纳入单位成功任务成本。第
\ref{ch:week09-optimization}~章据此比较云端调用与自建服务的固定投入，
附录~\ref{app:ai-infra-lab}~提供数值计算。收费规则和性能会随服务、地区及
合同变化，算例参数均为教学设定。

\section{图形表达与运营界面}
\subsection{图表的视觉编码}
\subjectindex{shijuebianma}{视觉编码}

数据可视化将数据字段映射到位置、长度、颜色等视觉属性。Wilkinson 的图形语法、Wickham 的分层图形语法和 Vega-Lite 的编码语法，都通过数据、变换、几何标记、尺度和坐标描述图表的构成。

\begin{definition}[视觉编码]
视觉编码是把数据字段映射到位置、长度、颜色、大小、形状、文本、透明度等视觉通道的过程。图表的解释质量取决于编码是否符合数据类型、比较任务和读者感知能力。
\end{definition}

\begin{table}[htbp]
\centering
\caption{管理分析中的常见图表任务}
\begin{tabular}{P{0.19\linewidth}P{0.33\linewidth}P{0.36\linewidth}}
\toprule
\rowcolor{BDMTableHead}
任务 & 推荐图表 & 适用场景\\
\midrule

时间趋势 & 折线图、柱线组合 & 收入、订单量、待处理订单随时间变化。\\
类别比较 & 排序条形图 & 区域、客户分层、产品类别比较。\\
分布形态 & 直方图、箱线图 & 订单金额、交付周期、客单价分布。\\
目标达成 & 子弹图、目标线 & 收入目标、准时率目标、库存上限。\\
异常定位 & 散点图、标注点 & 识别高收入低准时率对象。\\
结构占比 & 堆叠条形图 & 渠道结构、客户分层结构。\\
\bottomrule
\end{tabular}
\end{table}

Cleveland and McGill 的图形感知研究表明，不同视觉通道支持的判断精度不同。通常，沿共同尺度的位置比较比面积、角度或颜色深浅更容易准确判断。因此，管理看板中应谨慎使用饼图、仪表盘和过度装饰的图形；当需要比较大小时，排序条形图往往比面积图更可靠。

\subsection{从统计表到看板}
\subjectindex{Dashboard}{Dashboard（看板）}

看板是面向特定角色和任务的信息界面。Shneiderman 的可视化信息寻求原则可概括为：先看整体，再缩放和过滤，最后按需查看细节。经营看板的信息组织可分为整体概览、比较、解释、明细查询与行动：
\[
\text{Overview}
\rightarrow
\text{Compare}
\rightarrow
\text{Explain}
\rightarrow
\text{Drill down}
\rightarrow
\text{Act}.
\]

一个基础经营看板至少应包含：
\begin{enumerate}
  \item 顶部 KPI 卡片：显示总收入、准时率、平均延期、待处理订单和质量问题；
  \item 趋势视图：展示原始时间序列和移动平均；
  \item 比较视图：按区域、客户分层或产品类别排序；
  \item 分布或异常视图：显示极端值、分位数和异常标记；
  \item 质量提示：展示本次刷新是否通过关键质量检查；
  \item 下钻入口：能回到订单、客户、产品或质量失败记录。
\end{enumerate}

\begin{convention}[看板最小可解释性]
分析项目中的看板至少应标注数据范围、刷新时间、指标口径、分子分母、目标线或比较基准、质量状态和下钻路径。这些信息构成指标追溯和结果审计的依据。
\end{convention}

\subsection{看板类型}

看板的内容取决于使用者的职责和任务。战略型看板概括长期目标，分析型看板支持分组比较，运营型看板提供当前待处理事项。

\begin{table}[htbp]
\centering
\caption{看板类型}
\begin{tabular}{P{0.20\linewidth}P{0.33\linewidth}P{0.34\linewidth}}
\toprule
\rowcolor{BDMTableHead}
类型 & 主要问题 & 设计重点\\
\midrule
战略型 & 组织目标是否达成？ & 少量高层 KPI、目标线、趋势和预警。\\
分析型 & 为什么出现差异？ & 分组比较、下钻、过滤、分布和异常。\\
运营型 & 现在需要处理什么？ & 实时或近实时状态、任务列表、阈值告警。\\
\bottomrule
\end{tabular}
\end{table}

收入波动、准时率与质量问题的分析需要分组比较和明细查询，适合采用分析型看板。若将同一数据做成运营型看板，则应更突出待处理订单、质量问题日期和需要人工介入的任务；若做成战略型看板，则应更突出目标达成、累计收入和趋势方向。

\subsection{看板的运营节奏}

看板的维护需与业务周期协调，具体包括以下安排。
\begin{itemize}
  \item 什么时候刷新，刷新失败如何处理；
  \item 谁每天或每周查看，谁对异常负责；
  \item 哪些异常触发工单、会议或自动通知；
  \item 指标口径变化如何公告和版本化；
  \item 质量检查失败时是否阻断发布；
  \item 管理行动是否被记录并回到后续复盘。
\end{itemize}

看板的运行需要明确刷新频率、告警条件、明细查询和处理责任。指标异常出现后，负责人根据相关记录核查原因，记录处理结果并更新后续观察。

\section{制造企业的运营看板}

\subsection{Rolls-Royce 的发动机维护与备件管理}
\subjectindex{RollsRoyce}{Rolls-Royce}
Databricks 的制造业新闻稿介绍 Rolls-Royce Civil Aerospace 利用发动机及服务网络数据支持库存规划和备件位置安排的实践\cite{rr05}。该航空发动机企业的数字化材料也讨论了数据整合在制造和服务中的应用\cite{rrdigital05}。

备件分析看板可显示发动机服务风险与备件可得率、各地点需求与库存，以及候选调拨和采购方案。故障概率影响备件需求，运输时延与服务成本影响备件位置。第~\ref{ch:week07-prediction}~章讨论需求预测，第~\ref{ch:week09-optimization}~章讨论资源配置。

在制造场景中，设备预警需要关联维修备件、替代工站能力与到期订单。相应的评价指标包括诊断时间、维护等待时间、缺料次数和最终交付表现。

\subsection{半导体制造看板的分析次序}
第一层展示到期 OTD、最终良率、理想良品工时 OEE、日均总 WIP 与完成批次周期，同时显示分母、截止时间和质量状态。第二层用产品组合和工站拆解总体变化。第三层才列出需要调查的批次、停机原因与工艺路径。

例如 OEE 下降、WIP 上升、Cycle Time 拉长可能提示设备可用性、换线或需求波动等候选机制。第2章的制造分析要素可用于整理设备、工艺和需求资料，再逐项检验候选机制。若 OTD 尚未下降，可能因为缓冲库存暂时吸收了问题；若产量提升但 WIP 急升，可能是向下游转移拥堵。

\subsection{人工智能辅助图表解释}
多模态模型的图表解释应以原始聚合表、单位、指标定义和图形为依据，并区分观测事实、候选机制和验证方法。例如“第6天不良率为 \(10\%\)，超过教学控制界限”是事实；“可能来自设备或产品组合变化”是假设；“检查同日产品结构和维护日志”是下一步查询。

可以用自动检查核对回答中的数字是否来自表，趋势方向是否正确，是否把相关写成因果。视觉识别、数字一致性和经营解释应分别检验。任何自动发布的注释都应记录所用数据与口径版本。

\section{看板计算与经营分析}
\subsection{描述统计计算}

配套样例提供数据 \path{data/sample/daily_operations.csv}，粒度为一行一个经营日期，字段包括收入、订单数、交付订单数、准时订单数、平均延期天数、待处理订单数和质量问题数。配套代码分布为：
\begin{itemize}
  \item 统计案例：\path{src/bdm_decision/cases/week05_statistics_dashboard.py}；
  \item HTML 看板工具：\path{src/bdm_decision/app/dashboard.py}。
\end{itemize}

最小运行方式为：
\begin{lstlisting}[style=pythonstyle,caption={运行第5章统计摘要}]
import bdm_decision.cases.week05_statistics_dashboard as week05

rows = week05.load_daily_operations()
revenue = [row.revenue for row in rows]

print(week05.descriptive_summary(revenue))
print(week05.dashboard_kpis(rows))
\end{lstlisting}

HTML 看板的生成入口见上述界面模块，运行环境与版本记录要求见附录~\ref{app:reproduction}。

该数据的主要结果为：
\begin{table}[htbp]
\centering
\caption{第5章样例数据的核心摘要}
\begin{tabular}{P{0.36\linewidth}P{0.24\linewidth}P{0.28\linewidth}}
\toprule
\rowcolor{BDMTableHead}
指标 & 数值 & 管理解释\\
\midrule
总收入 & 71350 & 累计经营规模。\\
平均日收入 & 5945.83 & 日收入典型水平，但受高值影响。\\
收入中位数 & 5100 & 更稳健的典型日收入。\\
收入 IQR & 2875 & 中间 50\% 日期的收入跨度。\\
收入异常值 & 13500 & 需要下钻的高收入日期。\\
已交付订单准时率 & 0.5833 & 需与事先约定的目标比较，并说明分子分母。\\
平均延期天数 & 0.75 & 已交付订单的加权平均延期。\\
最新待处理订单数 & 10 & 当前运营负荷。\\
质量问题总数 & 6 & 看板必须显示的可信度提示。\\
\bottomrule
\end{tabular}
\end{table}

这些数字构成看板的第一层摘要。进一步分析应调查准时率与目标的差距、高收入日期对应的客户或产品，以及质量问题与延期的时间分布。调查需要结合第~\ref{ch:week02-enterprise-data-model}~章的业务对象、第~\ref{ch:week03-sql-kpi}~章的 SQL/KPI 和第~\ref{ch:week04-data-quality}~章的质量失败记录。

\subsection{由观察形成管理解释}

看板分析报告应说明主要观察、比较依据、相关记录与后续调查，基本内容如下。
\begin{table}[htbp]
\centering
\caption{看板管理解释模板}
\begin{tabular}{P{0.22\linewidth}P{0.64\linewidth}}
\toprule
\rowcolor{BDMTableHead}
环节 & 内容\\
\midrule

经营状态 & 当前总量、典型水平、目标达成和质量状态。\\
关键波动 & 哪些日期、对象或指标显著偏离基准。\\
下钻问题 & 需要回到哪些订单、客户、产品或质量失败记录。\\
行动建议 & 哪些行动可以立即执行，哪些需要进一步验证。\\
解释边界 & 哪些结论只是描述性发现，不能被当作因果结论。\\
复盘入口 & 后续如何检查行动是否改变指标。\\
\bottomrule
\end{tabular}
\end{table}

第5章样例中，本期收入存在一个高值日期，平均日收入高于中位数，说明收入分布受单日高值影响；准时交付率为 0.5833，且质量问题总数为 6，因此看板应提示履约和数据质量共同构成风险；下一步应下钻 2026 年 3 月 14 日的订单构成，并检查质量问题日期是否与延期日期重合。由此形成了关于数据质量与履约结果关系的待检验假设。

\subsection{看板结果的核查}

分析项目中的看板应满足以下最低要求：
\begin{enumerate}
  \item 每个 KPI 卡片都能追溯到字段、分子、分母和时间窗口；
  \item 每张图注明比较对象、统计期间和变量定义；
  \item 至少一个图表展示时间趋势，且保留原始观测；
  \item 至少一个图表支持比较、分布或异常定位；
  \item 明确显示数据日期范围和质量状态；
  \item 提供从汇总指标下钻到明细对象的路径；
  \item 管理解释中区分描述、预测、因果和行动建议。
\end{enumerate}

\section*{本章小结}
\addcontentsline{toc}{section}{本章小结}

描述统计刻画观测资料的位置、离散程度和分布形态。分组比较需要说明权重，时间序列分析需要区分原始观测与平滑结果。图形通过视觉编码呈现这些特征，经营看板则按使用任务组织统计量、比较基准和明细查询。

产品组合分解把总体变化分为组内变化与权重变化；控制图根据稳定过程的概率模型判断观测偏离；Little 定律和排队模型说明在制品、吞吐率与停留时间的联系。对人工智能服务，应分别记录首输出、完整请求和验收后的任务时间，并按负载估计并发与费用。这些分析为研究经营变化提供可检验的假设。随机实验估计管理干预效果的方法见第~\ref{ch:week06-ab-testing}~章。

\chapterreading{探索性分析的思想见 Tukey\cite{tukey1977}，图形感知的实验依据见 Cleveland
与 McGill\cite{cleveland1984}，图形语法见 Wilkinson 与 Wickham
\cite{wilkinson2005,wickham2010}。看板设计可参考 Few\cite{few2006}。
Little 定律是存量、流量与周期联系的基础\cite{little05}；Kingman 的重交通研究
提供理解变异与等待的进阶入口\cite{kingman05}。图形用于呈现分布，排队模型描述相应的运行机制；模型服务的负载分析见\cite{aiinfra2026}。
干预效果的研究设计见第~\ref{ch:week06-ab-testing}~章与第~\ref{ch:week08-causality}~章。}

\section*{思考与习题}
\addcontentsline{toc}{section}{思考与习题}

\begin{exercise}[结构与队列]
验证本章产品组合例子的总体良率，计算固定初期组合后的末期良率。再用 \(\mu=12,\lambda=9\) 和 \(\lambda=11\) 分别计算 \(M/M/1\) 系统时间，解释只考核利用率的风险。
\end{exercise}
\begin{exercise}[控制图的边界]
取 \(p_0=0.05\)，分别计算 \(n=100\) 与 \(n=1000\) 的控制界限。说明为何图上同样的 \(8\%\) 不良率可能有不同统计含义，并列出使二项独立假设失效的两个生产机制。
\end{exercise}

\begin{exercise}
\begin{enumerate}
  \item 对 \path{data/sample/daily_operations.csv} 中的收入计算均值、中位数、标准差、四分位数、IQR 和变异系数。
  \item 使用 IQR 规则识别收入异常日期，并说明为什么异常不等于因果解释。
  \item 设计一个准时交付率 KPI 卡片，说明同时展示哪些分母、目标线和质量提示。
  \item 为同一数据设计一个战略看板和一个运营看板，说明二者信息密度和交互方式的差别。
  \item 任选一个图表，写出其数据字段、视觉通道、聚合方式和管理解释。
  \item 为第5章样例写一页管理解释，明确描述性发现、下钻问题和行动建议。
\end{enumerate}
\end{exercise}
